% Demo: morpheme alignment (\GlossMorphAlign) together with phantom
% alignment ([phantomalign]) past Quine corners declared as marks.
% Builds with pdflatex, xelatex and lualatex.
\DocumentMetadata{lang=en, pdfversion=2.0, pdfstandard=ua-2,
  testphase={phase-III}}
\documentclass[11pt]{article}
\usepackage[a4paper,margin=2.5cm]{geometry}
\usepackage{amssymb}% \ulcorner, \urcorner
\usepackage{linguexx}

% Quine corners as single-token commands, so the phantom scan can see them
% (manual, sec. "Aligning past brackets and judgment marks").
\newcommand{\qq}{\ensuremath{\ulcorner}}
\newcommand{\qqr}{\ensuremath{\urcorner}}
\GlossPhantomChars{*?([<\#\%\qq}

\title{Linguexx demo}
\date{}

\begin{document}
\maketitle

\section{Morpheme alignment}

By default a gloss column is aligned by word only. With
\verb|\GlossMorphAlign| every morpheme, hyphen included, starts at the same
point in all segmented tiers.

\noindent Off (by word):

\ex. \gll ev-ler-imiz-de otur-uyor-uz \\
          house-\lpzg{pl}-\lpzg{1pl.poss}-\lpzg{loc} sit-\lpzg{prog}-\lpzg{1pl} \\
\glt `We are sitting in our houses.'

\noindent On (by morpheme):

{\GlossMorphAlign
\ex. \gll ev-ler-imiz-de otur-uyor-uz \\
          house-\lpzg{pl}-\lpzg{1pl.poss}-\lpzg{loc} sit-\lpzg{prog}-\lpzg{1pl} \\
\glt `We are sitting in our houses.'

}

\noindent A three-tier gloss whose first line is unsegmented (it is set
whole and does not take part), a clitic boundary (\texttt{=}), and a braced
multi-word gloss that counts as one morpheme:

{\GlossMorphAlign
\ex. \glll Sie sahen's gestern \\
           sie sah-en=es gestern \\
           they see.\lpzg{pst}-\lpzg{3pl}={it} yesterday \\
\glt `They saw it yesterday.'

\ex. \gll Schwieger-mütter-n \\
          {in-law}-mother.\lpzg{pl}-\lpzg{dat.pl} \\
\glt `mothers-in-law'

}

\section{Quine corners and phantom alignment}


\noindent Phantom alignment off:

\ex. \gll \qq{}ev-ler\qqr{} \qq{}kitap-lar\qqr{} \\
          house-\lpzg{pl} book-\lpzg{pl} \\

\noindent Phantom alignment on:

{\GlossPhantomAlign
\ex. \gll \qq{}ev-ler\qqr{} \qq{}kitap-lar\qqr{} \\
          house-\lpzg{pl} book-\lpzg{pl} \\

}

\section{Both at once}

With both switched on, the phantom pad is counted into the first morpheme
of the gloss rather than added on top of it. The gloss carries its own
bracket but no corner, so its bracket stays under the object's and the
corner's width is padded behind it: \textit{house} sits under \textit{ev},
and every morpheme boundary after it still lines up.

{\GlossPhantomAlign\GlossMorphAlign
\ex. \gll (\qq{}ev-ler-imiz-de\qqr{}) otur-uyor-uz \\
          (house-\lpzg{pl}-\lpzg{1pl.poss}-\lpzg{loc}) sit-\lpzg{prog}-\lpzg{1pl} \\
\glt `We are sitting in our houses.'

\ex. \gll [\qq{}kitap-lar\qqr{}] \qq{}gel-di-ler\qqr{} \\
          [book-\lpzg{pl}] come-\lpzg{pst}-\lpzg{3pl} \\
\glt `(The) books came.'

}

\section{Alternatives}
{\GlossPhantomAlign
  \ex. \altn{This}{John}{My car} is \altn[l]{truly}{really}{?very} great!
  \z.
}  

\end{document}

%%% Local Variables:
%%% mode: LaTeX
%%% TeX-master: t
%%% End:
